\documentclass[11pt]{article}
\usepackage[margin=1in]{geometry}
\usepackage{graphicx}
\usepackage{booktabs}
\usepackage{longtable}
\usepackage{array}
\usepackage{enumitem}
\usepackage{fancyvrb}
\usepackage{amsmath}
\usepackage[hidelinks]{hyperref}

\title{Point Taken: An Argument-Coaching Study\\
\large Design, Measures, and Preregistered Analysis Plan}
\author{Corey Zhang \and Aakriti Kumar \and Steven Franconeri}
\date{August 2026}

\begin{document}

\maketitle

\begin{abstract}
\noindent
We describe the design and preregistered analysis plan for a between-subjects
experiment testing whether a between-round, large language model (LLM) coaching
intervention improves the quality of participants' argumentation over repeated
practice. The study adapts the practice-plus-coaching paradigm of Kumar et al.
(arXiv:2603.15245) from empathic communication to disagreement and argument.
Participants complete three rounds of a structured argument game against a live,
good-faith LLM co-player; those in the treatment condition additionally receive
personalized coaching between rounds. Argument quality is scored on a six-dimension
rubric by an LLM judge, following the evaluation framework of Kumar et al.
(\emph{Nature Machine Intelligence}, 2025). This document specifies the design,
the measures, the exact model prompts, and the confirmatory analyses; it is
frozen prior to the full data collection.
\end{abstract}

\section{Introduction}

Point Taken is a structured discussion game in which participants construct,
place, and revise arguments in a reflective format rather than in the free-form
manner of ordinary online debate. The present study examines whether AI-delivered
coaching can improve the quality of a participant's argumentation across repeated
practice, without disrupting gameplay.

The design adapts an established practice-plus-coaching paradigm. Kumar et al.
(arXiv:2603.15245) show that practicing with an LLM partner and receiving
personalized feedback improves human empathic communication. We ask the analogous
question for disagreement: does between-round coaching improve how well people
argue? We hypothesize that participants who receive coaching will show larger
gains in argument quality across rounds than participants who practice without
coaching.

The document is organized as follows. Sections~\ref{sec:design}--\ref{sec:measures}
specify the design, procedure, and measures. Section~\ref{sec:judge} describes the
LLM-as-judge scoring and its validation. Section~\ref{sec:analysis} states the
preregistered analyses. Section~\ref{sec:results} is a placeholder for the pilot
results. The exact model prompts are reproduced verbatim in
Appendix~\ref{sec:prompts}, and the point-level coaching layer is documented in
Appendix~\ref{sec:pointlevel}.

\section{Why this study, and why now}
\label{sec:whynow}

% Corey's own words (assembled from his reasoning). Read and adjust so every line
% is one you can defend out loud.
AI is being used more and more in how people argue and reason online, and learning
how to interact with it well is becoming important. Part of why online arguing
goes wrong is that people cannot read body language or tone, so they misread each
other. Prior work has shown that practicing with an AI partner and getting
feedback can improve empathic communication, but whether an AI coach can help
people construct better arguments has not been tested. Arguing well is a different
skill: it is about building points from logic and evidence rather than emotion.
This study tests that directly. My expectation is that the coach will make
fallacies visible to people and help them structure their arguments better.

\section{Study design}
\label{sec:design}

The experiment uses a between-subjects design with two conditions, to which
participants are randomly assigned (50/50) on first entry:

\begin{itemize}[nosep]
  \item \textbf{Control:} three rounds of the argument game with no coaching.
  \item \textbf{Treatment (AICoach):} the same three rounds, with a between-round
  coaching interface.
\end{itemize}

A between-subjects design is used to prevent contamination of the control
condition by the coaching intervention. In each round, a participant argues
against a live LLM co-player rather than another participant; this permits
asynchronous, single-participant data collection on Prolific without real-time
matching. Each round presents one topic, with topic order randomized per
participant so that topic is not confounded with round.

The co-player is instructed to argue in good faith (a reasonable interlocutor
rather than an adversarial one). This choice isolates the construct of interest,
namely how well a participant argues against a fair partner; an adversarial
partner would instead index a participant's resistance to bad-faith tactics, a
distinct construct outside the scope of this study.

\begin{figure}[h]
\centering
\includegraphics[width=0.85\textwidth]{design_flow.pdf}
\caption{Experimental design. Participants are recruited on Prolific, complete a
baseline survey, and are randomly assigned (50/50) to the control or treatment
condition. Both conditions play three rounds against a good-faith LLM co-player;
the treatment condition additionally receives coaching between rounds. Each round
is scored by an LLM judge on six dimensions. The primary outcome is the change in
the composite score across rounds, compared between conditions.}
\label{fig:design}
\end{figure}

\section{Procedure: rounds and interaction requirements}
\label{sec:procedure}

Each round must contain sufficient interaction to (i) yield a scorable sample of
the participant's argumentation and (ii) provide the coach with material to
respond to, while remaining short enough to sustain participant engagement on
Prolific. Accordingly, a round is defined in terms of exchanges rather than
isolated contributions: a participant must contribute at least three argument
turns per round, at least one of which must directly respond to the co-player
rather than introduce an unrelated point. This requirement ensures mutual
engagement and renders the \emph{engagement} dimension measurable. Across three
rounds this yields approximately nine to twelve participant messages, or roughly
15--20 minutes of task time. The exact minimum thresholds are noted as an open
parameter in Section~\ref{sec:open}.

\section{Measures}
\label{sec:measures}

Each round is scored by an LLM judge on \textbf{six dimensions}, each on a
five-point scale ($1 =$ not at all present, $5 =$ strongly present). Following the
preregistered structure of Kumar et al. (arXiv:2603.15245), the dimensions divide
into three prescriptive (helpful) and three proscriptive (counterproductive)
behaviors:

\begin{itemize}[nosep]
  \item \textbf{Helpful (shown to treatment participants):} reasoning quality,
  engagement, constructive stance.
  \item \textbf{Counterproductive (analysis only, not shown):} fallacies,
  evasion, hostility.
\end{itemize}

Two round-level measures are derived:
\begin{align*}
  \texttt{composite} &= \text{mean of the three helpful dimensions} \quad (1\text{--}5),\\
  \texttt{analysis\_overall} &= \textstyle\sum(\text{helpful}) - \sum(\text{counterproductive}).
\end{align*}
The \texttt{composite} is the primary dependent measure; \texttt{analysis\_overall}
is the full six-dimension score, mirroring the prescriptive-minus-proscriptive
composite of the source paradigm.

Only the three helpful dimensions are displayed to treatment participants. This
is a deliberate design decision: displaying all six risks inducing participants to
optimize a single aggregate score, whereas the objective of the intervention is to
shift argumentation strategy. The three counterproductive dimensions correspond to
the point-level fallacy categories in Appendix~\ref{sec:pointlevel}
(overgeneralization to fallacies, irrelevant rebuttal to evasion, personal attack
and straw man to hostility), so that the fine-grained tile feedback and the
round-level rubric assess the same behaviors at two levels of granularity.

Participants additionally complete a baseline survey (topic stances and four
self-reported disagreement-style items) and an exit survey.

\subsection{Why these six dimensions}
\label{sec:whydims}

% Corey's own words (assembled from his reasoning), with the supporting citations.
% Read and adjust so every line is one you can defend out loud.
The six dimensions split into three that help an argument and three that hurt it.
On the helpful side: \emph{reasoning quality}, because good reasoning uses logic
and evidence while weak reasoning is just opinion --- this is the acceptability and
sufficiency of a claim in Wachsmuth et al.'s (2017) argument-quality taxonomy;
\emph{engagement}, meaning you respond to the other side's strongest point rather
than a weaker version (dialectical relevance, and the opposite of straw-manning);
and \emph{constructive stance}, where you build on and acknowledge their point
instead of ignoring its premise, which matches the respect and openness codes of
the Discourse Quality Index (Steenbergen et al., 2003). On the counterproductive
side: \emph{fallacies} --- I focus on the four most common and highest-impact ones
I saw in the earlier game data (vague claims, over-broad claims, causal leaps, and
irrelevant responses) rather than rarer textbook fallacies; \emph{evasion}, when
the other point is not acknowledged and no counterargument is even attempted; and
\emph{hostility}, attacking the person emotionally instead of the argument (an
ad hominem, again a failure of DQI respect). The three helpful dimensions capture
what good arguing adds; the three counterproductive ones capture what it should
avoid, mirroring the prescriptive/proscriptive split of the paradigm we replicate.

\section{LLM-as-judge and validation}
\label{sec:judge}

Round scoring is performed by a single LLM call per round
(\texttt{claude-sonnet-4-5}); the scoring prompt is reproduced in
Appendix~\ref{sec:prompts}. Because the primary outcome depends on this judge, we
validate it rather than assume its reliability, following the evaluation framework
of Kumar et al. (\emph{Nature Machine Intelligence}, 2025).

The judge is currently zero-shot. Expert annotations do not yet exist for this
task; we will therefore construct them. Following the pilot, an expert coder
(and, where feasible, a second coder to establish inter-rater reliability) will
score a random subsample of approximately 20--30 rounds on the same
six-dimension, five-point rubric. These expert scores serve two purposes: as
few-shot exemplars in the judge prompt, and as a reference against which to
compute agreement between the LLM and human coders (Krippendorff's $\alpha$),
benchmarked against human--human agreement on the overlapping subset. Dimensions
exhibiting weak agreement will have their anchors revised before the confirmatory
analyses are conducted. This validation is performed on pilot data, prior to the
full run.

\section{Preregistered analysis plan}
\label{sec:analysis}

The confirmatory analyses below are frozen prior to the full data collection. The
pilot ($n = 5$) serves only to validate the data pipeline and estimate task
duration, and is not used for inference.

\paragraph{Primary hypothesis (H1).} Treatment participants improve in argument
quality across rounds to a greater extent than control participants. This is
tested with a linear mixed-effects model,
\[
  \texttt{composite} \sim \text{condition} \times \text{round}_c + (1 \mid \text{participant}),
\]
where $\text{round}_c$ is round centered ($0,1,2$). The condition-by-round
interaction is the confirmatory test; a positive coefficient supports H1, and is
reported with its 95\% confidence interval. Should the mixed model fail to
converge at the available sample size, the preregistered fallback is a Welch
two-sample $t$-test on each participant's round-1-to-round-3 change score
(a difference-in-differences contrast).

\paragraph{Secondary hypotheses (H2--H3).} The same model is applied to
\texttt{analysis\_overall} (the full six-dimension measure) and to the summed
counterproductive score, testing whether coaching reduces counterproductive
behaviors across rounds relative to control.

\paragraph{Exploratory analyses (H4).} Per-dimension trajectories are examined
descriptively to identify which dimensions are most affected; these are not
corrected for multiple comparisons and are labeled exploratory.

\paragraph{Prespecified exclusions.} Participants are excluded if they do not
complete all three rounds; fail an effort check (e.g., empty or single-character
arguments across an entire round); submit a duplicate Prolific identifier; or have
any scored round that reverted to fallback text owing to a coach or opponent error
(flagged via the exported \texttt{model\_name} field). The number of exclusions and
reasons are reported.

\paragraph{Sample and inference.} The pilot ($n = 5$) is for pipeline validation
only. The full run targets approximately 30 participants per condition; at this
scale the study is treated as an estimate of the effect size and its confidence
interval rather than a fully powered confirmatory test, with a formal power
analysis for any subsequent confirmatory study based on the effect observed here.

All analyses are executed from \texttt{analysis/e18\_analysis.ipynb} on the
study's data export. The repository is organized to match the replication
repositories of the two source papers (a single analysis notebook together with a
data export and a written, preregistered plan; see Section~\ref{sec:refs}).

\section{Results}
\label{sec:results}

\textit{Placeholder; completed after data collection. All values are produced by
\texttt{analysis/e18\_analysis.ipynb} from the data export.}

\paragraph{Sample.} $N$ participants ($\_\_$ treatment, $\_\_$ control) following
$\_\_$ exclusions. Median completion time $\_\_$ minutes.

\paragraph{Primary result.} Mean composite score by round and condition:

\begin{center}
\begin{tabular}{lccc}
\toprule
Condition & Round 1 & Round 2 & Round 3 \\
\midrule
Control & \_\_ & \_\_ & \_\_ \\
Treatment & \_\_ & \_\_ & \_\_ \\
\bottomrule
\end{tabular}
\end{center}

Condition-by-round interaction: $b = \_\_$, 95\% CI $[\_\_,\,\_\_]$, $p = \_\_$.

\begin{figure}[h]
\centering
% \includegraphics[width=0.6\textwidth]{trajectory.png}
\caption{Composite argument-quality score by round and condition.}
\end{figure}

\paragraph{Secondary and validation.} Effects on \texttt{analysis\_overall} and
the counterproductive score: $\_\_$. LLM--human agreement (Krippendorff's
$\alpha$) per dimension, with the human--human benchmark: $\_\_$.

\section{Relationship to prior work}
\label{sec:refs}

The design and analysis follow two studies directly, and the repository is
structured to match their public replication materials.

\begin{itemize}[nosep]
  \item Kumar et al., \emph{Practicing with Language Models Cultivates Human
  Empathic Communication} (arXiv:2603.15245). Source of the practice-plus-coaching
  paradigm and the preregistered six-dimension (prescriptive minus proscriptive)
  scoring. Replication materials:
  \url{https://github.com/aakriti1kumar/replication-data-and-code-practicing-with-llms-cultivates-human-empathic-communication};
  Zenodo \url{https://doi.org/10.5281/zenodo.20703371}.
  \item Kumar et al., \emph{When Large Language Models are Reliable for Judging
  Empathic Communication} (\emph{Nature Machine Intelligence}, s42256-025-01169-6;
  arXiv:2506.10150). Basis for the LLM-as-judge validation procedure. Replication
  materials:
  \url{https://github.com/aakriti1kumar/replication-data-and-code-when-LLMs-reliable-empathic-communication};
  Zenodo \url{https://doi.org/10.5281/zenodo.17230987}.
\end{itemize}

\appendix

\section{Model prompts (verbatim)}
\label{sec:prompts}

The following are the exact prompts used in the live study.

\subsection{Round-scoring (judge) prompt}

Each dimension carries explicit 1/3/5 anchors, and rationales must cite specific
text. The reasoning-quality anchors operationalize the acceptability, relevance,
and sufficiency criteria of the argument-quality taxonomy (Wachsmuth et al.,
2017); the constructive-stance and hostility anchors follow the respect and
openness constructs of the Discourse Quality Index (Steenbergen et al., 2003).

\begin{Verbatim}[fontsize=\small]
You are an expert annotator scoring how well a participant argued in one
round of a structured debate game. Score each of six behaviors on a 1-5
scale using the anchors below. Judge only the participant's arguments, in
the context of the topic and the opponent's points, and base every score
on specific evidence in their text.

HELPFUL behaviors (constructive argumentation):
1. reasoning_quality — are the claims acceptable, relevant to the point,
   and sufficiently supported (a reason, example, mechanism, or evidence)?
   1 = bare assertions or non-sequiturs; 3 = partial or general support;
   5 = each key claim is specific and backed by a clear reason or example.
2. engagement — does the participant take on the opponent's strongest
   point and represent it fairly, rather than talking past it?
   1 = ignores the opponent; 3 = responds to a weaker version;
   5 = directly answers the opponent's strongest point, stated fairly.
3. constructive_stance — critiques the idea (not the person), grants valid
   points, stays open?
   1 = closed or dismissive; 3 = civil but no openness;
   5 = engages the idea and concedes or builds on a fair point.

COUNTERPRODUCTIVE behaviors (score HIGHER when MORE present):
4. fallacies — overgeneralization, false cause, exaggeration.
   1 = none; 3 = one clear instance; 5 = pervasive.
5. evasion — ignoring points, shifting subject, whataboutism.
   1 = stays on point; 3 = dodges once; 5 = repeatedly avoids.
6. hostility — personal attacks, straw-manning, dismissiveness.
   1 = none; 3 = one loaded remark; 5 = repeated bad-faith attacks.
{examples}
Topic: {topic}
The participant's arguments this round:
{arguments}

Score strictly against the anchors; each rationale must quote or point to
specific text. Respond with ONLY a JSON object with all six keys plus a
"feedback" field (one or two forward-looking tips for the next round).
\end{Verbatim}

The \texttt{\{examples\}} field receives expert-coded few-shot exemplars, following
the finding of Kumar et al. (\emph{Nature Machine Intelligence}, 2025) that a
detailed framework plus a small number of scored examples improves judge
reliability over zero-shot. Two illustrative examples currently seed this slot and
will be replaced with expert-coded rounds prior to the full run
(Section~\ref{sec:judge}).

\subsection{Between-round coaching prompt (primary intervention)}

The coaching shown to treatment participants between rounds uses a tutor-style
structure: identify a strength, focus on the participant's lowest-scoring helpful
dimension, and provide a concrete rewrite of one weaker argument. This mirrors the
personalized-feedback depth of the practice-plus-coaching paradigm
(arXiv:2603.15245).

\begin{Verbatim}[fontsize=\small]
You are an argument coach giving one participant feedback between rounds
of a debate game. You have the arguments they just made and their rubric
scores. Help them argue better next round, the way a good writing tutor
would: specific, warm, and focused on one or two changes.

Write feedback in three short parts, addressed to "you":
1. Strength: name one thing they did well, quoting their own words.
2. Focus: take the LOWEST-scoring helpful dimension and explain, in plain
   language, the single change that would move it most. Tie it to a
   specific argument of theirs.
3. Try this: rewrite ONE of their weaker arguments into a stronger version
   that keeps their side and intent, so they see the change concretely.

Under 120 words. Encouraging, no jargon, no numeric scores in the text.

Topic: {topic}
Their arguments: {arguments}
Helpful-dimension scores (1-5): reasoning_quality={rq},
engagement={en}, constructive_stance={cs}.

Return only the feedback text.
\end{Verbatim}

\subsection{Ask-the-coach prompt (optional follow-up)}

\begin{Verbatim}[fontsize=\small]
You are a supportive debate coach. The participant just argued a round.
Their arguments:
{arguments}

They ask: {question}

Answer in 2-3 sentences: specific, encouraging, actionable.
\end{Verbatim}

\subsection{Co-player prompt (content step)}

\begin{Verbatim}[fontsize=\small]
You are writing a single debate argument tile.
Topic: "{topic}"
You argue {side_label} the topic.
Responding to: "{target_text}"
Argument angle: {angle}

Rules:
- ONE grammatically complete sentence. Period at the end.
- Hard character limit; write a complete sentence that fits.
- No question marks. No fallacies.
- Argue in good faith: make the strongest FAIR point for your side.
  No cheap shots, no straw-manning, no exaggeration; a reasonable
  sparring partner, not a troll.
- Sound like a thoughtful student, not a textbook.
- Return your argument text only.
\end{Verbatim}

\section{Point-level coaching layer}
\label{sec:pointlevel}

Within a round, the treatment condition also provides point-level feedback on
individual argument tiles. This layer evaluates a single tile in context and
determines whether to surface a reasoning problem and how to express it in plain
language. It prioritizes precision over recall: a missed weak argument is
preferable to a false positive, since spurious feedback interrupts play and erodes
trust in the system.

The first-wave categories, with participant-facing labels in italics, are:
lack of specificity (\emph{Needs more detail}), overgeneralization
(\emph{Too broad}), false causation (\emph{Causal leap}), and irrelevant rebuttal
(\emph{Needs clearer connection}). Rarer categories (exaggeration, anecdotal
evidence, appeal to authority, straw man, whataboutism) are retained for research
analysis only, as they are less reliably detectable from a short tile.

Feedback generation proceeds in stages: evaluate the tile in context, determine
whether a target issue is present, determine whether it is severe enough to
surface, and, if so, produce a brief explanation and an optional rewrite. Any
rewrite must preserve the participant's original stance; a rewrite that alters the
side undermines the participant's intent.

\section{Open parameters}
\label{sec:open}

\begin{enumerate}[nosep]
  \item The exact minimum number of turns per round and the strictness of the
  requirement to reply directly to the co-player (Section~\ref{sec:procedure}).
  \item Final scoring anchors for each of the six dimensions, in particular
  constructive stance, which may exhibit a ceiling if it rewards only the absence
  of hostility.
  \item The composition of the validation coding team (a single expert coder, or
  two coders to establish inter-rater reliability).
  \item The judge model. The source evaluation framework used Gemini 2.5 Pro and
  Claude 3.7 Sonnet; the present study uses Claude Sonnet 4.5.
\end{enumerate}

\end{document}
